% ==================================================
% components_lcd.tex
% LCD icon and mode table components
% ==================================================

% ------------------------------
% Compact LCD icon table for raw LaTeX renderers
% Signatures:
%   \begin{HBLcdIconTable} ... \end{HBLcdIconTable}
%   \HBLcdIconRow{no}{image path or basename}{icon/name}{description}
% ------------------------------
\expandafter\providecommand\csname HBcomp_lcd_no_col_width\endcsname{5.2mm}
\expandafter\providecommand\csname HBcomp_lcd_icon_width\endcsname{8mm}
\expandafter\providecommand\csname HBcomp_lcd_icon_height\endcsname{8mm}
\expandafter\providecommand\csname HBcomp_lcd_icon_col_width\endcsname{11mm}
\expandafter\providecommand\csname HBcomp_lcd_label_col_width\endcsname{0.24\linewidth}
\expandafter\providecommand\csname HBcomp_lcd_table_tabcolsep\endcsname{2.0pt}
\expandafter\providecommand\csname HBcomp_lcd_table_row_stretch\endcsname{0.82}
\expandafter\providecommand\csname HBtype_lcd_no_font_size\endcsname{5.8pt}
\expandafter\providecommand\csname HBtype_lcd_no_font_leading\endcsname{6.4pt}
\expandafter\providecommand\csname HBtype_lcd_label_font_size\endcsname{6.2pt}
\expandafter\providecommand\csname HBtype_lcd_label_font_leading\endcsname{6.8pt}
\expandafter\providecommand\csname HBtype_lcd_body_font_size\endcsname{6.0pt}
\expandafter\providecommand\csname HBtype_lcd_body_font_leading\endcsname{6.8pt}

% Circled row numbers, matching the V2.0 JE-1000F master's LCD table.
\providecommand{\HBCircledNum}[1]{%
  \begin{tikzpicture}[baseline=(n.base)]
    \node[draw=\csname HBtype_body_color\endcsname, circle, inner sep=0.28mm,
          line width=0.28pt,
          font=\HBFontRegular\fontsize{\csname HBtype_lcd_no_font_size\endcsname}
                {\csname HBtype_lcd_no_font_leading\endcsname}\selectfont,
          text=\csname HBtype_body_color\endcsname] (n) {#1};
  \end{tikzpicture}%
}

\providecommand{\HBTypeLcdNo}[1]{%
  {\centering\HBCircledNum{#1}\par}%
}

\providecommand{\HBTypeLcdLabel}[1]{%
  {\HBFontSemiBold
   \color{\csname HBtype_body_color\endcsname}%
   \fontsize{\csname HBtype_lcd_label_font_size\endcsname}
            {\csname HBtype_lcd_label_font_leading\endcsname}\selectfont
   \RaggedRight #1}%
}

\providecommand{\HBTypeLcdText}[1]{%
  {\HBFontRegular
   \color{\csname HBtype_body_color\endcsname}%
   \fontsize{\csname HBtype_lcd_body_font_size\endcsname}
            {\csname HBtype_lcd_body_font_leading\endcsname}\selectfont
   \RaggedRight #1}%
}

\newtcolorbox{hblcdicontableframe}{
  enhanced,
  clip upper,
  colback=white,
  colframe=LineK40,
  arc=\csname HBcomp_table_outer_arc\endcsname,
  boxrule=\csname HBcomp_table_outer_rule\endcsname,
  boxsep=0mm,
  left=0mm,right=0mm,top=0mm,bottom=0mm,
}

% longtable instead of a framed tabular: the tcolorbox frame ("clip upper")
% cannot break across pages, so a tall icon table overflowed the 130x185
% text block, clipped its last rows, and collided with the page number
% (reports/typography_gap/2026-07-03 #1/#2). longtable breaks between rows;
% \endhead re-draws the top rule on every continuation page, and each row's
% own \hline closes the bottom edge before a break.
% Trade-off: the rounded outer corner (HBcomp_table_outer_arc) is dropped
% for this table.
\newenvironment{HBLcdIconTable}{%
  \begingroup
  \arrayrulecolor{LineK40}
  \arrayrulewidth=\dimexpr\csname HBcomp_table_outer_rule\endcsname/2\relax
  \ifdim\arrayrulewidth<0.2pt\arrayrulewidth=0.2pt\fi
  \setlength{\tabcolsep}{\csname HBcomp_lcd_table_tabcolsep\endcsname}
  \renewcommand{\arraystretch}{\csname HBcomp_lcd_table_row_stretch\endcsname}
  \setlength{\LTpre}{0pt}%
  \setlength{\LTpost}{0pt}%
  \setlength{\LTleft}{0pt}%
  \setlength{\LTright}{\fill}%
  \begin{longtable}{|
    >{\centering\arraybackslash}m{\csname HBcomp_lcd_no_col_width\endcsname}|
    >{\centering\arraybackslash}m{\csname HBcomp_lcd_icon_col_width\endcsname}|
    >{\RaggedRight\arraybackslash}m{\csname HBcomp_lcd_label_col_width\endcsname}|
    >{\RaggedRight\arraybackslash}m{\dimexpr\linewidth
      -\csname HBcomp_lcd_no_col_width\endcsname
      -\csname HBcomp_lcd_icon_col_width\endcsname
      -\csname HBcomp_lcd_label_col_width\endcsname
      -8\tabcolsep-5\arrayrulewidth\relax}|}
  \hline\endhead
}{%
  \end{longtable}
  \endgroup
  \par
}

\providecommand{\HBLcdIconRow}[4]{%
  \HBTypeLcdNo{#1}%
  &
  \HBImageOrPlaceholder
    {\csname HBcomp_lcd_icon_width\endcsname}
    {\csname HBcomp_lcd_icon_height\endcsname}
    {#2}%
  &
  \HBTypeLcdLabel{#3}%
  &
  \HBTypeLcdText{#4}%
  \tabularnewline\hline
}

% ------------------------------
% Compact LCD mode table for operation pages
% Signatures:
%   \begin{HBLcdModeTable}{image path or basename}
%   \HBLcdModeFirstGroup{mode}{action1}{text1}{action2}{text2}{action3}{text3}
%   \HBLcdModeSecondGroup{mode}{action1}{text1}{action2}{text2}{action3}{text3}
%   \end{HBLcdModeTable}
% ------------------------------
\expandafter\providecommand\csname HBcomp_lcd_mode_image_col_width\endcsname{22mm}
\expandafter\providecommand\csname HBcomp_lcd_mode_image_width\endcsname{20mm}
\expandafter\providecommand\csname HBcomp_lcd_mode_image_height\endcsname{22mm}
\expandafter\providecommand\csname HBcomp_lcd_mode_state_col_width\endcsname{21mm}
\expandafter\providecommand\csname HBcomp_lcd_mode_action_col_width\endcsname{15mm}
\expandafter\providecommand\csname HBcomp_lcd_mode_table_tabcolsep\endcsname{2.0pt}
\expandafter\providecommand\csname HBcomp_lcd_mode_table_row_stretch\endcsname{0.9}

\newtcolorbox{hblcdmodetableframe}{
  enhanced,
  colback=white,
  colframe=LineK40,
  arc=\csname HBcomp_table_outer_arc\endcsname,
  boxrule=\csname HBcomp_table_outer_rule\endcsname,
  boxsep=0mm,
  left=0mm,right=0mm,top=0mm,bottom=0mm,
}

% Rebuilt without \multirow (typography_gap/2026-07-03 #5, the 45.8pt
% Overfull \vbox): \multirow does not grow its rows, so a state label
% taller than the three action rows it spanned spilled out of the table
% and "clip upper" hid the damage. Structure is now nested tables —
% outer (image | groups) -> one row per group (state | actions) ->
% action rows (action | text) — so every row grows to fit its tallest
% cell and nothing can overflow. "clip upper" is also removed from the
% frame: a future overflow must be visible, never silently clipped.
% Macro signatures are unchanged.
\makeatletter
\newdimen\HB@lcdmoderestw
\newdimen\HB@lcdmodeactblockw
\newdimen\HB@lcdmodetextw
\newenvironment{HBLcdModeTable}[1]{%
  \par\noindent
  \begin{hblcdmodetableframe}
  \begingroup
  \arrayrulecolor{LineK40}
  \arrayrulewidth=\dimexpr\csname HBcomp_table_outer_rule\endcsname/2\relax
  \ifdim\arrayrulewidth<0.2pt\arrayrulewidth=0.2pt\fi
  \setlength{\tabcolsep}{\csname HBcomp_lcd_mode_table_tabcolsep\endcsname}
  \renewcommand{\arraystretch}{\csname HBcomp_lcd_mode_table_row_stretch\endcsname}
  % width of the group block (everything right of the image column)
  \HB@lcdmoderestw=\dimexpr\linewidth
    -\csname HBcomp_lcd_mode_image_col_width\endcsname
    -2\tabcolsep-3\arrayrulewidth\relax
  % width of the actions block (right of the state column)
  \HB@lcdmodeactblockw=\dimexpr\HB@lcdmoderestw
    -\csname HBcomp_lcd_mode_state_col_width\endcsname
    -2\tabcolsep-\arrayrulewidth\relax
  % width of the description column inside the actions block
  \HB@lcdmodetextw=\dimexpr\HB@lcdmodeactblockw
    -\csname HBcomp_lcd_mode_action_col_width\endcsname
    -4\tabcolsep-\arrayrulewidth\relax
  \begin{tabular}{|
    >{\centering\arraybackslash}m{\csname HBcomp_lcd_mode_image_col_width\endcsname}|
    @{}>{\centering\arraybackslash}m{\HB@lcdmoderestw}@{}|}
  \hline
  \HBImageOrPlaceholder
    {\csname HBcomp_lcd_mode_image_width\endcsname}
    {\csname HBcomp_lcd_mode_image_height\endcsname}
    {#1}%
  &
  \begin{tabular}[c]{
    >{\RaggedRight\arraybackslash}m{\csname HBcomp_lcd_mode_state_col_width\endcsname}|
    @{}>{\centering\arraybackslash}m{\HB@lcdmodeactblockw}@{}}
}{%
  \end{tabular}%
  \tabularnewline
  \hline
  \end{tabular}
  \endgroup
  \end{hblcdmodetableframe}
  \par
}

\newcommand{\HB@lcdmodeactions}[6]{%
  \begin{tabular}[c]{
    >{\RaggedRight\arraybackslash}m{\csname HBcomp_lcd_mode_action_col_width\endcsname}|
    >{\RaggedRight\arraybackslash}m{\HB@lcdmodetextw}}
  \HBTypeLcdLabel{#1} & \HBTypeLcdText{#2}\tabularnewline \hline
  \HBTypeLcdLabel{#3} & \HBTypeLcdText{#4}\tabularnewline \hline
  \HBTypeLcdLabel{#5} & \HBTypeLcdText{#6}%
  \end{tabular}%
}

\newcommand{\HBLcdModeFirstGroup}[7]{%
  \HBTypeLcdLabel{#1}
  &
  \HB@lcdmodeactions{#2}{#3}{#4}{#5}{#6}{#7}%
  \tabularnewline \hline
}

\newcommand{\HBLcdModeSecondGroup}[7]{%
  \HBTypeLcdLabel{#1}
  &
  \HB@lcdmodeactions{#2}{#3}{#4}{#5}{#6}{#7}%
}
\makeatother
